% !TeX root = ../main.tex

\chapter{Introducción}

\section{Objetivo}

El objetivo principal de este Trabajo de Fin de Grado es emplear y validar un modelo de aprendizaje profundo capaz de extrapolar la reverberación tardía de una Respuesta al Impulso de Sala (RIR) a partir de su tramo temprano de 50 ms \cite{Lin25}. El problema se formula como una inferencia condicionada, donde la cabeza de la RIR aporta la información estructural necesaria para estimar la dinámica temporal y espectral de la cola reverberante.

Desde el punto de vista técnico, el trabajo persigue tres metas concretas. En primer lugar, probar una arquitectura codificador-decodificador eficiente que respete hipótesis físicas del decaimiento acústico. En segundo lugar, definir un procedimiento de entrenamiento reproducible sobre bases de datos de RIR simuladas. En tercer lugar, evaluar el comportamiento del modelo con métricas espectro-temporales y métricas físicas, de forma que la validación no sea solo numérica, sino también acústicamente interpretable.

El alcance incluye la implementación del flujo completo de datos, la adaptación e implementación de la arquitectura DECOR y su análisis experimental. Quedan fuera del alcance la simulación geométrica 3D en tiempo real y la optimización a bajo nivel para el despliegue.

\section{Justificación}

La extrapolación de respuestas impulsionales de sala es un problema relevante porque generar bien la reverberación tardía influye directamente en la calidad final del audio y en la utilidad de las simulaciones acústicas. Los métodos geométricos clásicos, aunque están bien fundamentados físicamente, pueden ser costosos cuando se necesita producir muchas respuestas o trabajar con tiempos de cálculo ajustados.

En este contexto, un enfoque de aprendizaje profundo con sesgo físico ofrece una alternativa de interés: reduce el coste de inferencia y, al mismo tiempo, mantiene coherencia con el comportamiento acústico esperado del recinto. Esta combinación es especialmente valiosa para evaluar de forma sistemática la extrapolación de colas reverberantes y generar casos de prueba acústicos con menos coste computacional.

La justificación metodológica del trabajo se apoya en la hipótesis de que el tramo temprano de la RIR contiene información suficiente para condicionar la estimación de la cola tardía, hipótesis desarrollada en el marco teórico y contrastada mediante EDC, $T_{60}$, DRR y pérdidas espectrales multirresolución.

\section{Impacto}

La relevancia del proyecto se centra en la extrapolación eficiente de RIR para análisis acústico y generación de datos en entornos de investigación. Frente a enfoques clásicos de simulación, la propuesta busca una relación más favorable entre coherencia física y coste computacional.

El impacto esperado se articula en tres planos. En el plano técnico, se persigue reducir significativamente el tiempo de generación de colas reverberantes respecto a métodos puramente físicos. En el plano acústico, se busca mantener coherencia de decaimiento energético y distribución espectral mediante un decodificador con estructura inspirada en modelos de reverberación. Respecto a sus aplicaciones, el sistema puede emplearse para aumento de datos, diseño sonoro y experimentación rápida en entornos donde la obtención de RIR medidas es limitada o costosa.

\section{Requisitos}

Las especificaciones del sistema se formulan con criterios medibles y verificables, orientados a trazabilidad experimental y validación cuantitativa.

\begin{itemize}
    \item \textbf{Requisitos de Entrada/Salida de Señal}
    \begin{itemize}
        \item Procesamiento monoaural con frecuencia de muestreo estricta de 48 kHz.
        \item Ventana de entrada acotada a $N_{head}=2400$ muestras.
        \item Horizonte de predicción de cola fijado en $N_{tail}=45600$ muestras.
    \end{itemize}

    \item \textbf{Requisitos de Coherencia Acústica}
    \begin{itemize}
        \item Conservación de continuidad energética entre cabeza y cola sintetizada.
        \item Error acotado sobre la curva de decaimiento energético (EDC).
        \item Desajustes controlados en métricas físicas de referencia como $T_{60}$ y DRR.
        \item Consistencia espectral en múltiples resoluciones de STFT.
    \end{itemize}

    \item \textbf{Requisitos del Entrenamiento}
    \begin{itemize}
        \item El entrenamiento debe ser estable de principio a fin, sin errores que paren el proceso.
        \item Si repetimos el experimento en las mismas condiciones, los resultados deben salir prácticamente iguales.
        \item El sistema debe ir guardando automáticamente la mejor versión del modelo durante el entrenamiento.
    \end{itemize}
\end{itemize}

